\documentclass[11pt,a4paper]{article}
\usepackage[italian]{babel}
\usepackage[margin=2.5cm]{geometry}
\usepackage{parskip}
\usepackage{amsmath, amssymb, amsthm}
\usepackage{xcolor}
\usepackage{listings}
\usepackage{tikz}
\usepackage[most]{tcolorbox}
\usepackage{hyperref}

% --- Riquadri con bordo nero, senza numero ---
\tcbset{nota/.style={colback=white, colframe=black, boxrule=0.5pt,
  sharp corners}}
\NewTColorBox{definizione}{}{nota,
  before upper={\textbf{Definizione.}\ }}
\NewTColorBox{osservazione}{}{nota,
  before upper={\textbf{Osservazione.}\ }}
\theoremstyle{definition}
\newtheorem*{esempio}{Esempio}
\newtheorem*{esercizio}{Esercizio}

% --- Codice C ---
% Uso: \begin{codice} ... \end{codice}; \begin{risultato} ... \end{risultato}
%      \lstinline|printf("ciao");| per il codice dentro al testo
\lstdefinestyle{c}{
  language=C,
  basicstyle=\ttfamily\small,
  keywordstyle=\color{blue}\bfseries,
  commentstyle=\color{gray}\itshape,
  stringstyle=\color{red!70!black},
  numbers=left,
  numberstyle=\tiny\color{gray},
  numbersep=8pt,
  columns=flexible,
  keepspaces=true,
  breaklines=true,
  showstringspaces=false,
  tabsize=4,
  morekeywords={include, define},
  literate={à}{{\`a}}1 {è}{{\`e}}1 {é}{{\'e}}1 {ì}{{\`i}}1
           {ò}{{\`o}}1 {ù}{{\`u}}1
}
\lstset{style=c}
\newtcblisting{codice}{listing only, nota, left=6mm,
  listing options={style=c}}
\newtcblisting{risultato}{listing only, colback=black!5, colframe=black,
  boxrule=0.5pt, sharp corners,
  listing options={basicstyle=\ttfamily\small, numbers=none, language={}}}

% --- Diagrammi di flusso (simboli standard) ---
\usetikzlibrary{shapes.geometric, arrows.meta, positioning}
\tikzset{
  terminale/.style = {ellipse, draw, minimum width=2.5cm, minimum height=0.9cm},
  io/.style        = {trapezium, trapezium left angle=70, trapezium right angle=110,
                      draw, minimum width=2.5cm, minimum height=0.9cm},
  processo/.style  = {rectangle, draw, minimum width=2.5cm, minimum height=0.9cm},
  decisione/.style = {diamond, draw, aspect=2, minimum width=2.5cm},
  freccia/.style   = {thick, -{Stealth}}
}

\title{Esercizi sui diagrammi di flusso}
\author{Marco Cerbella}
\date{Lezione 2 -- 23 settembre 2026}

\begin{document}
\maketitle

\section{Somma di 10 e 20}

Si usa una variabile \texttt{sum} inizializzata a zero; si leggono i due numeri,
si sommano salvando il risultato in \texttt{sum} e infine si stampa.

\begin{enumerate}
    \item Inizializza $sum = 0$ (processo)
    \item Leggi i numeri (I/O)
    \item $sum = 10 + 20$ (processo)
    \item Stampa $sum$ (I/O)
\end{enumerate}

\begin{center}
\begin{tikzpicture}[node distance=0.7cm]
  \node (a) [terminale] {Inizio};
  \node (b) [processo, below=of a] {$sum = 0$};
  \node (c) [io, below=of b] {Leggi i numeri};
  \node (d) [processo, below=of c] {$sum = 10 + 20$};
  \node (e) [io, below=of d] {Stampa $sum$};
  \node (f) [terminale, below=of e] {Fine};
  \draw [freccia] (a) -- (b);
  \draw [freccia] (b) -- (c);
  \draw [freccia] (c) -- (d);
  \draw [freccia] (d) -- (e);
  \draw [freccia] (e) -- (f);
\end{tikzpicture}
\end{center}

\section{Somma di 5 numeri}

Servono due variabili: \texttt{sum} (il risultato) e \texttt{count} (quanti
numeri sono stati letti). Si usa un \emph{ciclo}: i passi si ripetono finché la
condizione è vera.

\begin{enumerate}
    \item $sum = 0$, $count = 0$ (processo)
    \item Leggi $n$ (I/O)
    \item $sum = sum + n$ e $count = count + 1$ (processo)
    \item $count < 5$? (decisione)
    \item Se sì, torna al passo 2; altrimenti stampa $sum$ (I/O)
\end{enumerate}

\begin{center}
\begin{tikzpicture}
  \node (a) at (0,0) [terminale] {Inizio};
  \node (b) at (0,-1.4) [processo] {$sum = 0,\ count = 0$};
  \node (c) at (0,-2.8) [io] {Leggi $n$};
  \node (d) at (0,-4.2) [processo] {$sum = sum + n$};
  \node (d2) at (0,-5.4) [processo] {$count = count + 1$};
  \node (e) at (0,-7.2) [decisione] {$count < 5$};
  \node (f) at (0,-9.2) [io] {Stampa $sum$};
  \node (g) at (0,-10.6) [terminale] {Fine};
  \draw [freccia] (a) -- (b);
  \draw [freccia] (b) -- (c);
  \draw [freccia] (c) -- (d);
  \draw [freccia] (d) -- (d2);
  \draw [freccia] (d2) -- (e);
  \draw [freccia] (e) -- node[right]{No} (f);
  \draw [freccia] (f) -- (g);
  \draw [freccia] (e.west) -- node[above]{Sì} ++(-2.3,0) |- (c.west);
\end{tikzpicture}
\end{center}

Il programma C corrispondente:
\begin{codice}
#include <stdio.h>

int main() {
    int sum = 0, count = 0, n;
    while (count < 5) {
        scanf("%d", &n);
        sum = sum + n;
        count = count + 1;
    }
    printf("sum = %d\n", sum);
    return 0;
}
\end{codice}

\section{Stampare ``Hello World'' 10 volte}

Anche qui si usa un ciclo con un contatore.

\begin{enumerate}
    \item $count = 0$ (processo)
    \item Stampa ``Hello World'' (I/O)
    \item $count = count + 1$ (processo)
    \item $count < 10$? Se sì torna al passo 2, altrimenti termina.
\end{enumerate}

\begin{center}
\begin{tikzpicture}
  \node (a) at (0,0) [terminale] {Inizio};
  \node (b) at (0,-1.4) [processo] {$count = 0$};
  \node (c) at (0,-2.8) [io] {Stampa ``Hello World''};
  \node (d) at (0,-4.2) [processo] {$count = count + 1$};
  \node (e) at (0,-6) [decisione] {$count < 10$};
  \node (f) at (0,-8) [terminale] {Fine};
  \draw [freccia] (a) -- (b);
  \draw [freccia] (b) -- (c);
  \draw [freccia] (c) -- (d);
  \draw [freccia] (d) -- (e);
  \draw [freccia] (e) -- node[right]{No} (f);
  \draw [freccia] (e.west) -- node[above]{Sì} ++(-2.6,0) |- (c.west);
\end{tikzpicture}
\end{center}

\begin{codice}
#include <stdio.h>

int main() {
    int count = 0;
    do {
        printf("Hello World\n");
        count = count + 1;
    } while (count < 10);
    return 0;
}
\end{codice}

\begin{osservazione}
Il diagramma esegue la stampa \emph{prima} di controllare la condizione, quindi
il costrutto C fedele è \lstinline|do ... while|. Con un \lstinline|while|
la condizione verrebbe controllata prima del corpo.
\end{osservazione}

\section{Login a un account (da risolvere liberamente)}

Esercizio aperto: non esiste una sola soluzione. Una possibile sequenza:
\begin{enumerate}
    \item Inserisci l'indirizzo del sito nel browser (I/O)
    \item Si carica la home page (processo)
    \item Inserisci email e password (I/O)
    \item Credenziali valide? (decisione)
    \item Se no: errore di login (processo) e torna al passo 3;
          se sì: mostra l'account (I/O) e termina.
\end{enumerate}

\begin{center}
\begin{tikzpicture}
  \node (a) at (0,0) [terminale] {Inizio};
  \node (b) at (0,-1.3) [io] {Inserisci URL};
  \node (c) at (0,-2.6) [processo] {Si carica la home};
  \node (d) at (0,-3.9) [io] {Inserisci email e password};
  \node (e) at (0,-5.8) [decisione] {Credenziali valide?};
  \node (f) at (0,-7.8) [io] {Mostra l'account};
  \node (g) at (0,-9.1) [terminale] {Fine};
  \node (err) at (4.8,-5.8) [processo] {Errore di login};
  \draw [freccia] (a) -- (b);
  \draw [freccia] (b) -- (c);
  \draw [freccia] (c) -- (d);
  \draw [freccia] (d) -- (e);
  \draw [freccia] (e) -- node[right]{Sì} (f);
  \draw [freccia] (f) -- (g);
  \draw [freccia] (e.east) -- node[above]{No} (err.west);
  \draw [freccia] (err.north) |- (d.east);
\end{tikzpicture}
\end{center}
\end{document}
